\bookchapter{第五章：等手機更懂你——基礎設定}

\chapterintro{花幾分鐘設定，體驗提升巨大。呢啲設定就好似畀手機「調音」，等佢更符合你嘅使用習慣。}

\booksection{一、聲音與震動}

設個鍾意嘅鈴聲同通知聲。靈活使用靜音模式（開會）、勿擾模式（瞓覺），世界瞬間清靜。

\begin{itemize}
  \item \textbf{鈴聲}：揀一個你聽到嘅、唔尷尬嘅。
  \item \textbf{震動模式}：靜音嗰陣靠震動提醒。
  \item \textbf{勿擾模式}：可以設定例外，比如允許通訊錄來電。
  \item \textbf{音量分離}：媒體音量、鈴聲音量、鬧鐘音量可以分開調。
\end{itemize}

\booksection{二、顯示與亮度}

\begin{itemize}
  \item \textbf{自動亮度}：等手機根據環境光自動調節，慳心。
  \item \textbf{護眼/夜覽模式}：過濾藍光，夜晚睇手機舒服啲。
  \item \textbf{深色模式}：慳電（對OLED螢幕特別有效），且夜間唔刺眼。好多App都支援深色模式。
  \item \textbf{螢幕刷新率}：部分安卓機可調，越高（如120Hz）滑動越絲滑。
  \item \textbf{自動鎖屏}：設定15秒或30秒，安全又慳電。
  \item \textbf{字體大小}：覺得字細？調大啲。
\end{itemize}

\booksection{三、桌布、主題與桌面}

換張好睇嘅桌布（靜態或動態），心情都變咗。安卓品牌手機常有主題商店，能整套換圖標同風格。

\textbf{桌面小組件}：天氣、日程、音樂控制放桌面，一步直達。iOS同安卓都支援，長按桌面空白處添加。

\booksection{四、效率設定：等手機成為你嘅「外掛大腦」}

呢部分係「一旦設定就返唔到轉頭」嘅功能，每個都可以等你嘅日常效率顯著提升。

\booksubsection{專注 / 免打擾模式：奪回被通知奪走嘅注意力}

\begin{itemize}
  \item \textbf{核心思想}：唔係「少睇手機」，而係「喺該專注嗰陣專注，該聯絡嗰陣再聯絡」。
  \item \textbf{使用建議}：\newline · 工作模式：僅允許「星標」聯絡人同必要 App 通過；\newline · 睡眠模式：自動開啟，僅允許鬧鐘；\newline · 健身 / 閲讀 / 會議：根據目前活動定製。
  \item \textbf{進階玩法}：用自動化 + 位置 / 時間觸發，等模式自動切換。例如：iOS 快捷指令可設定「到達公司即自動進入工作模式」。
\end{itemize}

\booksubsection{通知管理：等紅點唔再「奪命連環 call」}

\begin{itemize}
  \item \textbf{關閉營銷通知}：90\% 嘅推送其實你並唔需要。
  \item \textbf{按重要性排序}：將人際通訊類（電話、短訊、即時通訊）置頂；新聞、購物、社交平臺類可設為「靜默推送」。
  \item \textbf{分組通知}：將同類通知摺疊埋一齊，避免一次解鎖被幾十條消息淹沒。
\end{itemize}

\booksubsection{快捷手勢：等操作快人一步}

\begin{itemize}
  \item \textbf{截屏}：電源 + 音量下（iOS 全面屏）/電源 + 音量下（安卓通用）。
  \item \textbf{截長圖 / 滾動截屏}：普通截圖後點擊「長截圖」選項。
  \item \textbf{一鍵掃碼}：控制中心 / 快捷設定入面添加「掃一掃」圖標。
  \item \textbf{輕點背面（iPhone）}：喺 [設定] > [輔助功能] > [觸控] > [輕點背面] 中綁定「截屏」或「鎖屏」等動作。
  \item \textbf{語音喚醒}：Siri / 小愛同學 / 小藝 / 小布，嗌一聲就可以查天氣、設鬧鐘。
\end{itemize}

\booksubsection{輸入法：最被低估嘅效率神器}

\begin{itemize}
  \item \textbf{揀啱輸入法}：\newline · 搜狗 / 百度：詞庫大，流行語更新快；\newline · 訊飛：語音輸入識別率最高，開會記筆記神器；\newline · Gboard（谷歌）：簡潔、跨平臺同步強、支援滑行輸入；\newline · iOS 自帶：私隱最好，同系統深度集成。
  \item \textbf{必備開關}：智能糾錯、聯想詞、雲詞庫同步、上滑輸入數字符號。
  \item \textbf{私隱提示}：慎用「整句聯想」或「個性化推薦」功能，可能上傳輸入內容。
\end{itemize}

\booksubsection{桌面與多任務：等工作流順暢起嚟}

\begin{itemize}
  \item \textbf{App 分屏 / 畫中畫}：睇影片嗰陣回消息，上網嗰陣記筆記。
  \item \textbf{小窗 / 懸浮窗}：微信 / 抖音支援「小窗模式」，回覆完即收起。
  \item \textbf{常用 App 置頂}：Dock 欄（螢幕底部）淨係放最高頻嘅 4 個，減少查找時間。
\end{itemize}

\begin{tcbinfo}{「上工第一天」檢查清單}
攞到新手機或大版本更新後，建議花 10 分鐘過一遍呢啲設定：\newline (1) 鎖屏密碼 + 生物識別 (2) 通知分組 + 關閉營銷推送 (3) 字體大小 / 刷新率 / 深色模式\newline (4) 預設輸入法 / 瀏覽器 / 郵件 (5) 緊急聯絡人 / 醫療卡 (6) 備份設定

\end{tcbinfo}

\booksection{五、輔助功能與無障礙：科技向善嘅力量}

唔好被「輔助功能」呢個名嚇親，呢啲功能人人可以用。佢哋最初係為殘障人士設計嘅，但對普通人都好有用——臨時手受傷抱娃、地鐵上耳機壞咗聽唔清、戴老花鏡睇細字……都可以救場。講到底，無障礙設計體現咗\textbf{科技向善}嘅理念：等每個人都可以平等咁享受數碼生活。

\booksubsection{視障用戶嘅「眼睛」——螢幕閲讀器全攻略}

螢幕閲讀器係視障用戶使用手機嘅核心工具。佢會將螢幕上\textbf{所有可視內容朗讀出嚟}，用戶通過手勢同語音操作手機。下面分別介紹 iOS 「旁白」同安卓 TalkBack 嘅使用方法。

\begin{tcbinfo}{iOS 旁白（VoiceOver）}
\textbf{開啟路徑}：[設定] > [輔助功能] > [旁白]，打開開關。都可三擊側邊鍵（或主螢幕鍵）快速開關。

\textbf{基本手勢}：

\begin{enumerate}
  \item 單指滑動：逐元素瀏覽螢幕內容；
  \item 雙指輕掃：上下/左右整頁滾動；
  \item 單指雙擊：激活目前焦點元素（相當於「點擊」）；
  \item 三指輕掃：切換頁面；
  \item 單指三擊：放大或選擇（與普通三擊唔同）。
\end{enumerate}

\textbf{個性化設定}：[旁白] > [語音]，可調節\textbf{朗讀語速}（預設約 50\%，可調到 70-80\% 聽感更流暢）、\textbf{音調}、切換唔同語音包。iOS 自帶「增強語音」同「語音控制器」，前者能減少呼吸聲同停頓，聽感更自然。

\end{tcbinfo}

\begin{tcbinfo}{安卓 TalkBack}
\textbf{開啟路徑}：[設定] > [輔助功能] > [TalkBack]，打開開關。或長按音量上下鍵 3 秒快速啟用（部分品牌支援）。

\textbf{基本手勢}：

\begin{enumerate}
  \item 單指左右滑動：切換焦點元素；
  \item 單指雙擊：激活元素；
  \item 雙指上下/左右滑：滾動；
  \item 三指單擊：進入閲讀模式（連續朗讀）；
  \item 四指單擊：暫停/繼續朗讀。
\end{enumerate}

\textbf{個性化設定}：[TalkBack] > [文字轉語音設定]，可調整\textbf{語速}、\textbf{音調}、切換\textbf{語音引擎}（建議安裝 Google 文字轉語音引擎，朗讀更自然）。

\end{tcbinfo}

\booksubsection{等細字變大、等模糊變清晰——視覺調整}

\booksubsection{放手勢}

\begin{itemize}
  \item \textbf{iOS}：[設定] > [輔助功能] > [縮放]，開啟後\textbf{三指輕點螢幕}即可放大。預設全屏縮放，可選窗口縮放（局部放大窗）。\textbf{雙指開合可調節倍數（1.5x 至 15x）}，三指雙擊鎖定目前縮放比例。
  \item \textbf{安卓}：[設定] > [輔助功能] > [放大]，三指點擊放大，再三指點擊退出；或直接\textbf{三指點擊螢幕三下}啟用臨時放大（最多 8 倍）。
\end{itemize}

\booksubsection{高對比度文字 / 粗體文字}

\begin{itemize}
  \item \textbf{iOS}：[設定] > [輔助功能] > [顯示與文字大小]，開啟「粗體文字」、「增強對比度」、「降低透明度」等開關。「智能反轉」可喺淺色背景下等文字保持黑色，避免白底白字。
  \item \textbf{安卓}：[設定] > [顯示] > [字體大小]，可調節「細、普通、大、超大」四檔。華為/小米額外提供「粗體」同「高對比度文本」開關。
\end{itemize}

\booksubsection{色彩濾鏡（色盲友好）}

\begin{itemize}
  \item \textbf{iOS}：[設定] > [輔助功能] > [顯示與文字大小] > [色彩濾鏡]，提供\textbf{紅/綠/藍/黃四種色盲模式}，以及「灰度」同「反轉顏色」選項。建議先做一次色覺測試（系統自帶）再選對應濾鏡。
  \item \textbf{安卓}：[設定] > [輔助功能] > [顏色校正]，含\textbf{三種色盲模式}+ 「反轉顏色」開關。
\end{itemize}

\booksubsection{AI 圖像描述：等 AI 當你嘅「眼睛」}

當螢幕閲讀器遇到圖片嗰陣，傳統做法係跳過或淨係讀出「圖片」兩個字。但而家 AI 可以\textbf{用一句話描述圖片內容}，對識別驗證碼、讀圖、識別菜單好有幫助。

\begin{tcbinfo}{各系統 AI 圖像描述支援情況}
\begin{itemize}
  \item \textbf{iOS 15+}：喺 [設定] > [輔助功能] > [旁白] > [語音描述] 中開啟「圖像描述」。系統會調用\textbf{本地 AI 模型}自動生成圖片文字描述（需 iPhone 12 及以上）。
  \item \textbf{新版 TalkBack 應用}：Google TalkBack 喺較新版本中集成\textbf{圖片問答}功能，可以問「呢張圖入面有乜嘢」，AI 實時回答（唔嚴格依賴安卓系統版本，更新 TalkBack 應用即可體驗）。
  \item \textbf{第三方 App}：「Seeing AI」（iOS）、「Lookout」（安卓）、「小藝幫幫」（華為）都係優秀嘅 AI 視覺助手，可識別物品、讀檔案、識別人臉表情。
\end{itemize}

\end{tcbinfo}

\begin{tcbpractice}{使用場景}
試下對住一張街拍相片問「呢張圖入面有幾多人？佢哋喺度做乜嘢？」，或者打開外賣 App，等 AI 讀出菜單圖入面嘅菜品名——\textbf{呢個係真正嘅「科技等生活更平等」}。

\end{tcbpractice}

\booksubsection{聽障用戶嘅「耳朵」——實時字幕全攻略}

實時字幕（Live Captions）係\textbf{系統級嘅語音轉文字}功能，聽障用戶喺嘈雜環境、圖書館、外語場景下都特別有用。

\begin{tcbinfo}{實時字幕開啟與優化}
\textbf{iOS}：[設定] > [輔助功能] > [實時字幕]，打開開關。系統會檢測設備所有音頻（通話、影片、播客、FaceTime、語音備忘錄），自動顯示字幕喺螢幕上方。

\textbf{安卓 12+}：按\textbf{音量鍵}點擊彈出框中嘅「實時字幕」圖標；或喺 [設定] > [輔助功能] > [實時字幕] 中開啟。

\textbf{準確率優化技巧}：

\begin{enumerate}
  \item 盡量喺\textbf{安靜環境}下使用，遠離冷氣、機器轟鳴；
  \item 說話者\textbf{咬字清晰、語速適中}嗰陣識別率最高；
  \item 多人對話會降低準確率，可使用「淨係顯示一人」模式；
  \item 支援中文普通話、粵語、英語等主流語種，\textbf{小語種識別率較低}；
  \item 說話中途有背景音樂或人聲幹擾，識別會斷斷續續。
\end{enumerate}

\end{tcbinfo}

\booksubsection{通話場景雙通道：語音轉文字 + 文字轉語音}

當用戶\textbf{既聽唔到都講唔到嘢}嗰陣（如突發失聲、喉部手術恢復期、嚴重口吃），手機仲可以當「傳話員」。

\begin{tcbinfo}{通話實時轉譯（TTY / RTT）}
\textbf{iOS}：[設定] > [輔助功能] > [RTT/TTY]，開啟後通話嗰陣螢幕上會顯示鍵盤，對方嘅語音會被實時打成文字，\textbf{你打嘅字都會被念出嚟}畀對方聽。

\textbf{安卓}：[設定] > [輔助功能] > [實時轉錄] + [聲音通知]。「實時轉錄」將對方說話實時顯示為文字；「聲音通知」喺不便語音嗰陣將你嘅輸入轉為語音播放。

\textbf{兼容性}：

\begin{itemize}
  \item \textbf{RTT（Real-Time Text）}：雙方都需支援，文字邊打邊顯示，無延遲，主流電訊商已支援；
  \item \textbf{TTY（電傳打字）}：需專用 TTY 設備，已逐步被 RTT 取代；
  \item 微信/QQ 嘅\textbf{語音轉文字功能}都係替代方案：長按語音消息 「轉文字」。
\end{itemize}

\end{tcbinfo}

\booksubsection{面向所有人嘅「細但實用」功能}

\begin{itemize}
  \item \textbf{輕點背面}（iOS 14+）：[設定] > [輔助功能] > [觸控] > [輕點背面]，可綁定\textbf{雙擊/三擊}觸發截圖、鎖屏、控制中心等動作。雙手攞快遞嗰陣特別好用。
  \item \textbf{輔助觸控 / 懸浮球}：iOS [輔助功能] > [觸控] > [輔助觸控]，安卓 [輔助功能] > [懸浮導航]。虛擬按鈕替代手勢或物理按鍵，按鍵唔靈或單手操作嗰陣救場。
  \item \textbf{切換控制}：通過外部開關、單次點擊、頭部動作控制手機，\textbf{為肢體障礙用戶設計}。iOS [輔助功能] > [切換控制]，安卓 [輔助功能] > [開關控制]。
  \item \textbf{聲音識別}（iOS）：[輔助功能] > [聲音識別]，開啟後手機會「聽」門鈴、嬰兒喊聲、火警警報、警笛等聲音，\textbf{識別到就震動+螢幕通知}，對聽障用戶同睡眠期間好有價值。
  \item \textbf{LED 閃爍提醒}：iOS [輔助功能] > [音頻/視覺] > [LED 閃爍以示提醒]，安卓 [設定] > [通知] > [閃光燈提醒]。來電同通知嗰陣閃光燈閃爍，\textbf{嘈雜環境無聲都可以感知}。
  \item \textbf{單聲道音頻}：[設定] > [輔助功能] > [音頻] > [單聲道音頻]，將左右聲道合併。單耳聽力受損嘅用戶唔會丟失另一半聲音資訊。
\end{itemize}

\begin{tcbinfo}{蘋果嘅無障礙設計哲學}
蘋果喺無障礙領域投入巨大，每年都會更新同新增輔助功能。蘋果官網有專門嘅無障礙頁面，詳細介紹各項功能嘅使用方法。iOS 嘅無障礙功能通常比安卓更完善、體驗更流暢，呢個都係好多殘障用戶選擇蘋果嘅重要原因。

\end{tcbinfo}

\begin{tcbpractice}{花30秒體驗一下}
揾一段空閒時間，做以下體驗：

\begin{enumerate}
  \item 開啟\textbf{旁白 / TalkBack}，試下唔靠眼用三指滑動瀏覽微信；
  \item 打開\textbf{實時字幕}，對住手機說話睇識別準確率；
  \item 用\textbf{三指點擊}放大一段文字到 200\%，體驗視障用戶嘅日常；
  \item 關掉之前所有功能，諗下「如果我永遠需要呢啲功能，我會希望佢點樣設計」。
\end{enumerate}

\end{tcbpractice}
